% TOP_PROBLEM: {"id": 2975, "title": "Kirby Problem 4.99", "category_id": 7, "category": {"id": 7, "name": "topology", "display_name": "Topology", "description": "Properties preserved under continuous deformations.", "slug": "topology", "order_index": 7, "created_at": "2026-07-31T15:26:25.671Z"}, "status": "open"}
% FIRST_POSTED: null
% ENTRY_KIND: statement_only
% Imported from: batch13.tex; UnsolvedMath ID 2975
\documentclass[11pt]{article}
\usepackage[margin=25mm]{geometry}
\usepackage{fontspec}
\setmainfont{Latin Modern Roman}
\usepackage{amsmath,amssymb,mathtools}
\usepackage{unicode-math}
\setmathfont{Latin Modern Math}
\usepackage{xurl,hyperref}
\hypersetup{colorlinks=true,urlcolor=blue}
\setcounter{secnumdepth}{0}
\setlength{\parindent}{0pt}
\setlength{\parskip}{5pt}
\setlength{\emergencystretch}{3em}
\newcommand{\sref}[2]{\hyperlink{source:#1:#2}{[S#2]}}
\newcommand{\cataloguescope}[1]{\paragraph{Recorded target.}#1}
\begin{document}
% UnsolvedMath ID 2975; source catalogue code KP-4.99
\setcounter{section}{2974}
\section{Kirby Problem 4.99}\label{2975}
{\small\sffamily UnsolvedMath ID 2975 · Topology}
\subsection{Definitions and mathematical statement}

\paragraph{Shared notation.}$\mathbb N=\{1,2,\ldots\}$, and $\mathbb Z,\mathbb Q,\mathbb R,\mathbb C$ denote the integers, rationals, reals and complex numbers. Any other conventions are specified locally.


Let $(X,\omega)$ be a closed symplectic four-manifold and let $T\subset X$ be a smoothly embedded symplectic torus with $T\cdot T=0$. Symplectic means that $\omega$ is a closed nondegenerate two-form and $\omega|_T$ is an area form. 

For a smooth knot $K\subset S^3$, let $M_K$ be zero-framed Dehn surgery on $K$, using the longitude bounding a Seifert surface in its exterior. Let $m$ be a meridian of $K$, regarded as a framed circle in $M_K$, and set $T_m=m\times S^1\subset M_K\times S^1$. If $T\subset X$ is a torus with trivial normal bundle, Fintushel--Stern knot surgery is
\[
X_K=\bigl(X\setminus\operatorname{int}(T\times D^2)\bigr)
\cup_\varphi\bigl((M_K\times S^1)\setminus\operatorname{int}(T_m\times D^2)\bigr).
\]
Here $\varphi$ is the usual orientation-reversing identification of the boundary three-tori, preserving the normal-circle class $[\{\mathrm{pt}\}\times\partial D^2]$ in the fiber-sum convention. Choose the framing and boundary identification as part of the knot-surgery data. Equivalently, the replacement piece is $S^1$ times the knot exterior, with its Seifert-longitude class identified with the normal circle of $T$. \sref{2975}{3}

A knot $K$ is fibered if its exterior $S^3\setminus\operatorname{int}\nu K$ is a smooth fiber bundle over $S^1$ with fiber a compact connected Seifert surface, and the fibration sends a meridian to a degree-one loop in the base.
\paragraph{Statement recorded in the supplied source.}
For such $(X,T)$ and a knot $K$, if the knot-surgery manifold $X_K$ admits a symplectic structure, must $K$ be fibered?
\[
 X_K\text{ admits a symplectic form with its inherited orientation}\quad\Longrightarrow\quad K\text{ is fibered}.
\]
The orientation of $X_K$ is inherited from the orientation of $(X,\omega)$ and the fiber-sum gluing; the resulting symplectic form must induce that orientation. It need not agree with $\omega$ on the unchanged region. No simple-connectivity, complement-group, or cusp-embedding hypothesis is imposed. \sref{2975}{2}
\subsection{Short English statement}
If knot surgery on a symplectic torus yields a manifold that admits any symplectic structure, must the surgery knot be fibered?
\subsection{Sources}
\begin{description}
\item[\hypertarget{source:2975:1}{S1}] ULAM.AI and the UnsolvedMath Contributors, UnsolvedMath: A Curated Collection of Open Mathematics Problems, record 2975 (KP-4.99). CC BY 4.0. \url{https://huggingface.co/datasets/ulamai/UnsolvedMath}.
\item[\hypertarget{source:2975:2}{S2}] R. İnanç Baykur, Robion C. Kirby and Daniel Ruberman, K3—A New Problem List in Low-Dimensional Topology, Mathematical Surveys and Monographs 295 (2026), author version, Problem 4.99 and Remarks, PDF page 272. \url{https://math.berkeley.edu/sites/default/files/surv-295-ruberman-watermarked-author-pdf.pdf}.
\item[\hypertarget{source:2975:3}{S3}] Ronald Fintushel and Ronald J. Stern, Knots, Links, and 4-Manifolds, Inventiones Mathematicae 134 (1998), 363--400; arXiv:dg-ga/9612014v2, Introduction, pages 3--4. \url{https://arxiv.org/pdf/dg-ga/9612014v2}.
\end{description}
\label{end:2975}
\end{document}
